\section{Four Facets of the Round Trip}
\label{sec:facets}

The single condition $\Dmap\circ\Cmap=\mathrm{id}$ has four faces. Each facet is
a way the round trip can be examined, stressed, or made to fail.

\paragraph{1. Predictive power (temporal direction).}
Predictive power is the round trip executed in the temporal direction: expanding
from the compressed state forward into future phenomena. The error is
controlled and nonzero---a perfect identity on the past is neither expected nor
required; what matters is that the expansion's error stays within a known bound.
In the source corpus this is the face used to grade models: a model earns its
keep by expanding into testable, falsifiable predictions.

\paragraph{2. Perturbation (how the round trip is achieved).}
Perturbation theory is the \emph{implementation} of the round trip: represent
the decompression as a leading term plus successive corrections, and the round
trip closes within a convergence radius ($\Dmap\circ\Cmap\approx\mathrm{id}$).
Necessity enters as a rule: higher-order terms are added only when the leading
terms fail to close the round trip. This turns ``how much detail to keep'' into
a measurable question (which correction term first exceeds tolerance), not a
taste judgment.

\paragraph{3. Fitness (the landscape of thinking).}
There is no binary ``wrong''; there is a continuous \emph{fitness} measure on
candidate models. Progress is gradient descent on this landscape: incremental,
local moves are robust to noise, which is why stepwise refinement outperforms
occasional wholesale replacement. Fitness is the facet that makes the framework
compatible with empirical work: it is the quantity that experiments can measure.

\paragraph{4. Leading-term anchoring (robustness to randomness).}
Robustness to randomness comes from anchoring on the leading term: the leading
term occupies the widest valley in the fitness landscape, so it is the safest
thing to grasp first; perturbations are then handled term by term. ``Anchor the
leading term, then process corrections'' is the operational form of
perturbative thinking, and it explains why a model that misses its leading term
fails even when its details are elaborate.

Together the four facets partition the work of closing the round trip:
\emph{where} to expand (prediction), \emph{how} to expand (perturbation),
\emph{what good looks like} (fitness), and \emph{what to hold onto}
(leading-term anchoring).
